\section{Two-point evaluation from rational connectedness}
\label{app:variation}

We prove \cref{lem:variation} using only polynomial presentations
of maps from \(\PP^1\), analytic subsets, and the holomorphic
constant-rank theorem. This isolates the precise two-point variation
property used in the proof. For broader background on rationally
connected varieties, see \cite{KMM92,Kol99}.

We first record an elementary covering observation.
If \(A\) is a proper closed analytic subset of a connected
second-countable complex manifold \(M\), then \(A\) is covered
by countably many complex submanifolds of dimension smaller
than \(\dim M\). The covering submanifolds need not be contained
in \(A\).
To see this, choose a countable cover of \(A\) by sufficiently small connected
coordinate neighborhoods on which \(A\) is defined by holomorphic
functions. In each neighborhood choose a non-identically-zero defining
function \(f\). Such a function exists because a proper analytic subset
has empty interior.
At a zero of \(f\), let \(r\ge1\) be its finite order.
Some derivative of order \(r-1\) vanishes there and has
nonzero differential. Thus that point belongs to the regular
zero locus of one of the partial derivatives of \(f\).
There are countably many such derivatives and charts.
Their regular zero loci, with connected components separated,
give the asserted covering.

\begin{proof}[Proof of \cref{lem:variation}]
Embed \(X\) in \(\PP^N\). Every holomorphic map
\(C=\PP^1\to\PP^N\) is represented by a tuple of homogeneous
polynomials of a degree \(d\ge0\), with no common zero.
Indeed its pullback of \(\OO_{\PP^N}(1)\) is \(\OO_C(d)\),
and its homogeneous coordinate sections supply such a tuple.
For fixed \(d\), these tuples form an open subset \(M_d\)
of a finite-dimensional vector space. Openness follows
from compactness of \(C\).
There is a universal holomorphic evaluation map
\(C\times M_d\to\PP^N\).

We cover the tuples whose maps land in \(X\) by countably
many manifold families all of whose maps land in \(X\).
Start with the connected components of \(M_d\).
For a current connected submanifold \(M\), either every tuple
maps \(C\) into \(X\), in which case accept this family, or
there are \(\tau\in M\) and \(c\in C\) with
\(G(c,\tau)\notin X\). In the latter case
\[
\{\eta\in M:G(c,\eta)\in X\}
\]
is a proper closed analytic subset of \(M\), and contains
every desired tuple in \(M\).
Use the covering observation to replace it by countably
many lower-dimensional submanifolds, and repeat the same
test on each. All desired tuples survive. Extra points
introduced by the covering cause no problem: each new
manifold is tested anew. Dimension strictly decreases
whenever a family is not accepted. In dimension zero the
test accepts precisely the desired points. Finite depth
and countable branching give countably many accepted
families covering all the desired tuples.
On an accepted family the universal evaluation factors
holomorphically through the embedded submanifold \(X\).

Perform this construction for every \(d\). The resulting
countably many two-point evaluation maps to \(X\times X\)
have images containing general pairs: a rational curve
through two distinct points can be reparametrized to take
those values at the fixed distinct points \(a,b\).
By the convention in the introduction, these images contain a nonempty
Zariski open subset of \(X\times X\). It is also open in the usual
topology and has positive measure in smooth local coordinates.

Suppose, for a contradiction, that every one of these
evaluation maps has complex rank everywhere less than
\(2\dim X\). The image of each is then measure zero.
Here is a direct verification of that fact for a holomorphic
map \(\phi:M\to X\times X\) with the stated rank bound.
We may take \(M\) connected. On the open set of maximum
rank, the constant-rank theorem places the image locally
in a submanifold of smaller dimension than \(X\times X\).
Countably many local charts give a measure-zero image.
If the maximum rank is zero, this already treats all of \(M\).
Otherwise its complement is a proper closed analytic
subset. Cover that subset by countably many
lower-dimensional submanifolds and apply induction on
\(\dim M\). The strict rank bound persists on every such
restriction, including the extra points of a covering
submanifold. The zero-dimensional case is immediate.
This proves the asserted measure-zero property.

A countable union of measure-zero images cannot contain
that open subset. Some family therefore has a point where
two-point evaluation is submersive. Restrict its parameter
space to the nonempty open set where the differential is
surjective; call the resulting manifold \(T\).

Finally, a parameter tangent vector differentiates the
holomorphic family to a global section of \(g^*T_X\).
Surjectivity of the two-point evaluation differential
provides, for every \(w\in T_{g(b)}X\), such a section with
values \(0,w\) at \(a,b\). Projecting with \(P_i\) gives the
last assertion of \cref{lem:variation}.
\end{proof}

